\documentclass[10pt,a4paper]{article}

% ----- Encodage & langue -----
\usepackage[T1]{fontenc}
\usepackage[utf8]{inputenc}
\usepackage[french]{babel}
\usepackage{ProfCollege}

% ----- Mise en page -----
\usepackage{geometry}
\geometry{top=1.5cm,bottom=1.5cm,left=1.5cm,right=1.5cm}

% Colonnes
\usepackage{multicol}
\setlength{\columnsep}{0.75cm}            % espace entre colonnes
\setlength{\columnseprule}{0.2pt}      % épaisseur de la ligne
\def\columnseprulecolor{\color{Couleur8}} % couleur de la ligne

% Maths et figures
\usepackage{amsmath, amssymb}
\usepackage{tikz}
\usetikzlibrary{babel,arrows,calc,fit,patterns,plotmarks,shapes.geometric,shapes.misc,shapes.symbols,shapes.arrows,shapes.callouts, shapes.multipart, shapes.gates.logic.US,shapes.gates.logic.IEC, er, automata,backgrounds,chains,topaths,trees,petri,mindmap,matrix, calendar, folding,fadings,through,positioning,scopes,decorations.fractals,decorations.shapes,decorations.text,decorations.pathmorphing,decorations.pathreplacing,decorations.footprints,decorations.markings,shadows}

\usetikzlibrary{calc}
\usepackage{graphicx}
% Listes compactes
\usepackage{enumitem}
\setlist{nosep, itemsep=0.4em, parsep=0pt}
\setlist[enumerate,1]{label=\alph*.}

% Couleurs
\usepackage{xcolor}
\definecolor{exoblue}{HTML}{005F73}

%----------------------------------------------------------------------------------------
%	 COULEURS
%----------------------------------------------------------------------------------------

\definecolor{Couleur1}{HTML}{8b1f30}
\definecolor{Couleur2}{HTML}{725E74}
\definecolor{Couleur3}{HTML}{78985b}
\definecolor{Couleur4}{HTML}{db3e56}
\definecolor{Couleur5}{HTML}{487C4B}
\definecolor{Couleur6}{HTML}{CF6876}
\definecolor{Couleur7}{HTML}{A5C03F} %Vert
\definecolor{Couleur8}{HTML}{D36740} %Orange
\definecolor{Couleur9}{HTML}{7BB48A} %Bleu-Vert
\definecolor{Couleur10}{HTML}{EA8556} %Orange
\definecolor{Couleur11}{HTML}{E6B459} %Jaune
\definecolor{Couleur12}{HTML}{9AC8EB} %Bleu

\definecolor{Bord1}{HTML}{302635}
\definecolor{Bord2}{HTML}{725E74}
\definecolor{Bord3}{HTML}{938999}
\definecolor{Bord4}{HTML}{817888}
\definecolor{Bord5}{HTML}{487C4B}
\definecolor{Bord6}{HTML}{8C436B}
\definecolor{Bord7}{HTML}{A5C03F} %Vert
\definecolor{Bord8}{HTML}{D36740} %Orange

% ----- Police sans serif -----
\renewcommand{\familydefault}{\sfdefault}

% ----- Exercice -----
\newcounter{exo}
\newcommand{\exo}[1][]{%
  \refstepcounter{exo}%
  \textbf{\textcolor{Couleur8}{\\[0.3cm]\theexo.}}%
  \if\relax\detokenize{#1}\relax\else\ \textbf{#1}\fi\ %
}

\newcommand{\ligne}{\noindent\dotfill\vspace{0.6em}}



\begin{document}
\begin{Large}\textbf{\textcolor{Couleur8}{EXERCICES - Chapitre 17 - Durées}} \end{Large}
\begin{multicols}{2}

%1
% @see : https://coopmaths.fr/alea?uuid=4f8f4&id=6M4C-1&n=2&d=10&s=1&alea=d9Pr&cd=0&cols=1
\exo \textit{Comme l'exemple 1}

\begin{enumerate} 
	\item Convertir $237$ minutes en heures ($\text{h}$) et minutes ($\text{min}$).
	
	\ligne
	
	\ligne

	\ligne

	\ligne
	
	\ligne

	\ligne


	\item Convertir $148$ minutes en heures ($\text{h}$) et minutes ($\text{min}$).
	
	\ligne
	
	\ligne

	\ligne

	\ligne

	\ligne

	\ligne


\end{enumerate}

\exo \textit{Comme l'exemple 2} % @see : https://coopmaths.fr/alea?uuid=8b0f9&id=6M4C&n=2&d=10&s=3&alea=nG2p&cd=1&cols=1

\begin{enumerate} 
	\item Convertir $3\,256~\text{s}$ en heures, minutes et secondes.
	
	\ligne
	
	\ligne

	\ligne

	\ligne

	\ligne

	\ligne


	\item Convertir $9\,554~\text{s}$ en heures, minutes et secondes.
	
	\ligne
	
	\ligne

	\ligne

	\ligne

	\ligne

	\ligne


\end{enumerate}

\exo \textit{Comme l'exemple 3}% @see : https://coopmaths.fr/alea?uuid=8b0f9&id=6M4C&n=2&d=10&s=1&alea=lnJM&cd=1&cols=1

\begin{enumerate} 
	\item Convertir $1~\text{h}~59~\text{min}$ en secondes.
	
	\ligne
	
	\ligne

	\ligne

	\ligne


	\item Convertir $8~\text{h}$ en minutes.
	
	\ligne
	
	\ligne

	\ligne


\end{enumerate}

\exo % @see : https://coopmaths.fr/alea?uuid=8b0f9&id=6M4C&n=5&d=10&s=5&alea=5NzZ&cd=1&cols=1


\begin{enumerate} 
	\item Convertir $116~\text{h}$ en jours et heures.
	
	\ligne
	
	\ligne

	\ligne

	\ligne


	\item Convertir $542~\text{h}$  en semaines, jours et heures.
	
	\ligne
	
	\ligne

	\ligne

	\ligne
	
	\ligne

	\ligne



	\item Convertir $6~\text{h}$ en minutes.
	
	\ligne
	
	\ligne

	\ligne

	\ligne


	\item Convertir $6\,921~\text{s}$ en heures, minutes et secondes.
	
	\ligne
	
	\ligne

	\ligne

	\ligne

	\ligne

	\ligne



	\item Convertir $53~\text{min}$ en secondes.
	
	\ligne
	
	\ligne

	\ligne



\end{enumerate}

\exo% @see : https://coopmaths.fr/alea?uuid=185f8&id=6M4C-3&n=2&d=10&s=1&s2=false&s3=true&alea=Zs5b&cd=1&cols=1

\begin{enumerate} 
	\item Convertir $2\,$heures\,$24\,$minutes  en heures.
	
	\ligne
	
	\ligne

	\ligne

	\ligne



	\item Convertir $3\,$heures\,$42\,$minutes  en heures.
	
	\ligne
	
	\ligne

	\ligne

	\ligne




\end{enumerate}

\exo% @see : https://coopmaths.fr/alea?uuid=185f8&id=6M4C-3&n=2&d=10&s=3&s2=false&s3=true&alea=Eyf6&cd=1&cols=1


\begin{enumerate} 
	\item Convertir $216$ minutes en heures.
	
	\ligne
	
	\ligne

	\ligne

	\ligne



	\item Convertir $234$ minutes en heures.
	
	\ligne
	
	\ligne

	\ligne

	\ligne




\end{enumerate}

% @see : https://coopmaths.fr/alea?uuid=185f8&id=6M4C-3&n=2&d=10&s=3&s2=false&s3=true&alea=Eyf6&cd=1&cols=1

\begin{enumerate} 
	\item Convertir $216$ minutes en heures.
	
	\ligne
	
	\ligne

	\ligne

	\ligne



	\item Convertir $234$ minutes en heures.
	
	\ligne
	
	\ligne

	\ligne

	\ligne



\end{enumerate}

 % @see : https://coopmaths.fr/alea?uuid=6b3e4&id=6M4C-2&n=2&d=10&alea=zJwo&cd=1&cols=1
\exo Écrire les durées suivantes en heures et minutes.

\begin{enumerate} 
	\item $9,4~\text{h}=$
	
	\ligne
	
	\ligne

	\ligne

	\ligne


	\item $4,2~\text{h}=$
	
	\ligne
	
	\ligne

	\ligne

	\ligne


\end{enumerate}

\exo % @see : https://coopmaths.fr/alea?uuid=8013e&id=6M4C-4&n=4&d=10&s=3&s2=true&s3=true&alea=3nhH&cd=1&cols=1

\begin{enumerate} 
	\item Convertir $174$ minutes en format heures-minutes.
	
	\ligne
	
	\ligne

	\ligne

	\ligne



	\item Convertir $2{,}75$ heures en format heures-minutes.
	
	\ligne
	
	\ligne

	\ligne

	\ligne



	\item Convertir $192$ minutes en format heures-minutes.
	
	\ligne
	
	\ligne

	\ligne

	\ligne



	\item Convertir $1{,}75$ heures en format heures-minutes.
	
	\ligne
	
	\ligne

	\ligne

	\ligne


\end{enumerate}

\end{multicols}

\newpage

\begin{Large}\textbf{\textcolor{Couleur8}{CORRECTION - Chapitre 15 - Durées}} \end{Large}
\setcounter{exo}{0}
\exo

\begin{enumerate} 
	\item $237 = (3 \times 60) + 57$ donc $237$ minutes = ${\color[HTML]{F15929}\boldsymbol{3\,\textbf{h}\,57\,\textbf{min}}}$.
	\item $148 = (2 \times 60) + 28$ donc $148$ minutes = ${\color[HTML]{F15929}\boldsymbol{2\,\textbf{h}\,28\,\textbf{min}}}$.
\end{enumerate}
 

\exo

\begin{enumerate} 
	\item \opidiv[resultstyle={\color[HTML]{F15929}},lineheight=\baselineskip,columnwidth=2ex,voperation=top]{3256}{60}$3\,256~\text{s} = (54\times60~\text{s})+16~\text{s}={\color[HTML]{F15929}\boldsymbol{54~\textbf{min}~16~\textbf{s}}}$
	\item \opidiv[resultstyle={\color[HTML]{F15929}},lineheight=\baselineskip,columnwidth=2ex,voperation=top]{9554}{3600}\opidiv[resultstyle={\color[HTML]{F15929}},lineheight=\baselineskip,columnwidth=2ex,voperation=top]{2354}{60}\\$9\,554~\text{s} = (2\times3\,600~\text{s})+2\,354~\text{s} =2~\text{h}+(39\times60~\text{s})+14~\text{s}= {\color[HTML]{F15929}\boldsymbol{2~\textbf{h}~39~\textbf{min}~14~\textbf{s}}}$
\end{enumerate}
 

\exo

\begin{enumerate} 
	\item $1~\text{h}~59~\text{min} = (1\times3~600~\text{s}) + (59\times60~\text{s}) = 3\,600~\text{s}+3\,540~\text{s} = {\color[HTML]{F15929}\boldsymbol{7\,140~\textbf{s}}}$
	\item $8~\text{h} = 8\times60~\text{min} = {\color[HTML]{F15929}\boldsymbol{480~\textbf{min}}}$
\end{enumerate}
 

\exo

\begin{enumerate} 
	\item \opidiv[resultstyle={\color[HTML]{F15929}},lineheight=\baselineskip,columnwidth=2ex,voperation=top]{116}{24}\\$116~\text{h} = (4\times24~\text{h}) + 20~\text{h} = {\color[HTML]{F15929}\boldsymbol{4~\textbf{j}~20~\textbf{h}}}$
	\item \opidiv[resultstyle={\color[HTML]{F15929}},lineheight=\baselineskip,columnwidth=2ex,voperation=top]{542}{24}\opidiv[resultstyle={\color[HTML]{F15929}},lineheight=\baselineskip,columnwidth=2ex,voperation=top]{22}{7}\\$542~\text{h} = (22\times24~\text{h}) + 14~\text{h} = 22~\text{j}~14~\text{h} = (3\times7~\text{j}) + 1~\text{j} + 14~\text{h} = {\color[HTML]{F15929}\boldsymbol{3~\textbf{semaines}~1~\textbf{j}~14~\textbf{h}}}$
	\item $6~\text{h} = 6\times60~\text{min} = {\color[HTML]{F15929}\boldsymbol{360~\textbf{min}}}$
	\item \opidiv[resultstyle={\color[HTML]{F15929}},lineheight=\baselineskip,columnwidth=2ex,voperation=top]{6921}{3600}\opidiv[resultstyle={\color[HTML]{F15929}},lineheight=\baselineskip,columnwidth=2ex,voperation=top]{3321}{60}\\$6\,921~\text{s} = (1\times3\,600~\text{s})+3\,321~\text{s} =1~\text{h}+(55\times60~\text{s})+21~\text{s}= {\color[HTML]{F15929}\boldsymbol{1~\textbf{h}~55~\textbf{min}~21~\textbf{s}}}$
	\item $53~\text{min} = 53\times60~\text{s} = {\color[HTML]{F15929}\boldsymbol{3\,180~\textbf{s}}}$
\end{enumerate}
 

\exo

\begin{enumerate} 
	\item Pour convertir un temps au format heures-minutes en heures, on convertit les minutes en fractions d'heures. \\ Le nombre d'heures ne change pas : $2\text{\,h}$. \\On divise le nombre de minutes par $60$ pour le convertir en heures : $24\text{\,min} = \dfrac{24}{60}\text{\,h}=0{,}4\text{\,h}$ \\On somme : $2 + 0{,}4 = 2{,}4$.\\Ainsi, $2\text{\,h } 24\text{\,min } = {\color[HTML]{F15929}\boldsymbol{2{,}4}}\text{\,h}$.
	\item Pour convertir un temps au format heures-minutes en heures, on convertit les minutes en fractions d'heures. \\ Le nombre d'heures ne change pas : $3\text{\,h}$. \\On divise le nombre de minutes par $60$ pour le convertir en heures : $42\text{\,min} = \dfrac{42}{60}\text{\,h}=0{,}7\text{\,h}$ \\On somme : $3 + 0{,}7 = 3{,}7$.\\Ainsi, $3\text{\,h } 42\text{\,min } = {\color[HTML]{F15929}\boldsymbol{3{,}7}}\text{\,h}$.
\end{enumerate}
 

\exo

\begin{enumerate} 
	\item Pour convertir des minutes en heures, on divise par $60$ (car $60\text{\,min} = 1\text{\,h}$).

 
$216\text{\,min} = 216 \div 60 \text{\,h}= 3{,}6\text{\,h}$.\\Ainsi, $216\text{\,min} = {\color[HTML]{F15929}\boldsymbol{3{,}6}}\text{\,h}$.
	\item Pour convertir des minutes en heures, on divise par $60$ (car $60\text{\,min} = 1\text{\,h}$).

 
$234\text{\,min} = 234 \div 60 \text{\,h}= 3{,}9\text{\,h}$.\\Ainsi, $234\text{\,min} = {\color[HTML]{F15929}\boldsymbol{3{,}9}}\text{\,h}$.
\end{enumerate}
 

\exo

\begin{enumerate} 
	\item $9,4~\text{h}=9~\text{h}+\dfrac{4}{10}~\text{h}=9~\text{h}+(4\times6~\text{min})=$\,${\color[HTML]{F15929}\boldsymbol{9\,\textbf{h}\,24\,\textbf{min}}}$
	\item $4,2~\text{h}=4~\text{h}+\dfrac{2}{10}~\text{h}=4~\text{h}+(2\times6~\text{min})=$\,${\color[HTML]{F15929}\boldsymbol{4\,\textbf{h}\,12\,\textbf{min}}}$
\end{enumerate}
 

\exo

\begin{enumerate} 
	\item Pour convertir des minutes au format heures-minutes-secondes, on divise par $60$ pour obtenir le nombre d'heures, puis on traite le reste.

 
Le nombre d'heures entières correspond à la partie entière du quotient $174\text{\,min} \div 60$, soit $2\text{\,h}$. 

 
Il reste à trouver les minutes : $174-(2\times 60) =54\text{\,min}$.\\Ainsi, $174\text{\,min} = {\color[HTML]{F15929}\boldsymbol{2}}\text{\,h } {\color[HTML]{F15929}\boldsymbol{54}}\text{\,min}$.
	\item Pour convertir des heures au format heures-minutes, on sépare la partie entière (heures) et on convertit la partie décimale.\\$2{,}75\text{\,h} = 2\text{\,h} + 0{,}75\text{\,h}$

 
Conversion de la partie décimale en minutes :\\$0{,}75\text{\,h} = 0{,}75 \times 60\text{\,min} = 45\text{\,min}$.\\Ainsi, $2{,}75\text{\,h} = {\color[HTML]{F15929}\boldsymbol{2}}\text{\,h } {\color[HTML]{F15929}\boldsymbol{45}}\text{\,min}$.
	\item Pour convertir des minutes au format heures-minutes-secondes, on divise par $60$ pour obtenir le nombre d'heures, puis on traite le reste.

 
Le nombre d'heures entières correspond à la partie entière du quotient $192\text{\,min} \div 60$, soit $3\text{\,h}$. 

 
Il reste à trouver les minutes : $192-(3\times 60) =12\text{\,min}$.\\Ainsi, $192\text{\,min} = {\color[HTML]{F15929}\boldsymbol{3}}\text{\,h } {\color[HTML]{F15929}\boldsymbol{12}}\text{\,min}$.
	\item Pour convertir des heures au format heures-minutes, on sépare la partie entière (heures) et on convertit la partie décimale.\\$1{,}75\text{\,h} = 1\text{\,h} + 0{,}75\text{\,h}$

 
Conversion de la partie décimale en minutes :\\$0{,}75\text{\,h} = 0{,}75 \times 60\text{\,min} = 45\text{\,min}$.\\Ainsi, $1{,}75\text{\,h} = {\color[HTML]{F15929}\boldsymbol{1}}\text{\,h } {\color[HTML]{F15929}\boldsymbol{45}}\text{\,min}$.
\end{enumerate}


\end{document}